\section{Experimental Results}
\label{sec:exp-results}
The experimental section is organized around three complementary evaluation
views: answer-verifiable reasoning tasks, non-verifiable or proof-oriented
tasks, and official olympiad competition problems. 
% This structure separates
% standard automatic evaluation from settings that require proof-level judgment
% and from the competition-style benchmarks that most directly reflect the final
% goal of gold-medal-level reasoning.

\subsection{Benchmarks}

We organize evaluation into three benchmark families. The first family contains
answer-verifiable reasoning tasks, where correctness can be checked by a final
answer or a high-confidence automatic verifier. It includes AMO-Bench~\citep{an2025amobench},
AIME 2025 and AIME 2026\footnote{AIME official competition page:
\url{https://maa.org/math-competitions/aime}.}, AnswerBench from the
IMO-Bench evaluation suite~\citep{luong2025imobench}, and FrontierScience-Olympiad,
the Olympiad subset of FrontierScience~\citep{wang2026frontierscience}. These
benchmarks mainly test whether the model can produce correct final answers
under single-pass or fixed-budget inference.

The second family contains non-verifiable or proof-oriented tasks. We include
ProofBench from IMO-Bench~\citep{luong2025imobench}, which emphasizes proof
quality rather than only final-answer matching, and FrontierScience-Research, the
research subset of FrontierScience~\citep{wang2026frontierscience}. These tasks
are used to probe whether training improves rigorous reasoning and scientific
problem solving beyond answer-checkable settings.

The third family contains official olympiad competition problems, including
IMO 2025\footnote{International Mathematical Olympiad official archive:
\url{https://www.imo-official.org/}.}, USAMO 2026\footnote{USA
Mathematical Olympiad:
\url{https://maa.org/math-competitions/usamo}.} and IPhO (2024, 2025)\footnote{International Physics
Olympiad official site: \url{https://www.ipho-new.org/}.}. 
% We separate these tasks from
% general benchmark suites because they are closest to the final target setting:
% long-horizon, high-stakes mathematical or scientific reasoning with strict
% competition-style scoring. 
Detailed grading and verifier settings are summarized
in \cref{app:evaluation-details}.

\definecolor{MathReasoningBg}{RGB}{255,244,204}
\definecolor{ScienceReasoningBg}{RGB}{234,242,255}
\begin{table*}[t]
\centering
\caption{Performance on answer-verifiable reasoning tasks. Results for AnswerBench, AMO-Bench, AIME 25/26, and FrontierScience-Olympiad are averaged over 4, 8, 8, and 4 runs, respectively. FrontierScience-Olympiad abbreviates the Olympiad subset of FrontierScience. Avg. is the mean of AnswerBench, AMO-Bench, AIME 2025, AIME 2026, and FrontierScience-Olympiad. Within each comparison block, \textbf{bold} marks the best score and underline marks the second best.}
\label{tab:verifiable-single-pass}
\vspace{-5 pt}
\setlength{\tabcolsep}{2.2pt}
\renewcommand{\arraystretch}{1.14}
\small
\resizebox{\textwidth}{!}{%
\begin{tabular}{lcccccccc}
\toprule
\multirow{2}{*}{\textbf{Model}} & \multirow{2}{*}{\textbf{AnswerBench}} & \multirow{2}{*}{\textbf{AMO-Bench}} & \multirow{2}{*}{\textbf{AIME 25/26}} & \multicolumn{4}{c}{\textbf{FrontierScience-Olympiad}} & \multirow{2}{*}{\textbf{Avg.}} \\
\cmidrule(lr){5-8}
 & & & & \textbf{Physics} & \textbf{Chemistry} & \textbf{Biology} & \textbf{Overall} & \\
\midrule
% \multicolumn{5}{l}{\textit{Larger models}} \\
% \midrule
% % MiniMax-M2.7  & -- & -- & -- & -- \\
% % Kimi K2.6 & -- & -- & -- & -- \\
% DeepSeek-V3.2 & -- & -- & 91.7 / -- & -- \\
% % GLM-5.1 & -- & -- & -- & -- \\
% GPT-5.5 & -- & -- & 100.0 / -- & -- \\
% Gemini 3.1 Pro Thinking & -- & 65.5 & 96.7 / -- & -- \\
% \midrule
% \multicolumn{5}{l}{\textit{Similar-size models}} \\
% \midrule
P1-30B-A3B & 69.3\% & 41.3\% & 90.4\% / 89.6\% & 57.5\% & 57.5\% & \underline{27.5\%} & 54.5\% & 69.0\% \\
GLM-4.7-Flash & 73.8\% & 53.8\% & 91.3\% / 88.3\% & 54.5\% & 60.0\% & 17.5\% & 53.0\% & 72.0\% \\
Nemotron-Cascade-2 & \textbf{80.5\%} & 40.8\% & \underline{94.2\%} / 90.0\% & 56.0\% & 56.3\% & \textbf{30.0\%} & 53.5\% & 71.8\% \\
Qwen3.6-35B-A3B & \underline{78.0\%} & \underline{58.8\%} & 92.5\% / \underline{92.9\%} & \underline{65.5\%} & \textbf{74.4\%} & 25.0\% & \textbf{65.0\%} & \textbf{77.4\%} \\
Gemma-4-31B & 74.0\% & 39.3\% & 88.8\% / 91.3\% & \textbf{69.0\%} & 61.9\% & 17.5\% & 61.0\% & 70.9\% \\
% \midrule
% \multicolumn{5}{l}{\textit{Ours}} \\
% SFT & 59.8 & 31.0 & 81.3 / 85.8 & 55.0 \\
% SFT + Coarse RL & 77.3 & 57.3 & 92.9 / 90.0 & 60.3 \\
\rowcolor{ScienceReasoningBg!35}
\textbf{\textcolor{iclrdeepblue}{SU-01}} & 77.5\% & \textbf{59.8\%} & \textbf{94.6\%} / \textbf{93.3\%} & 62.5\% & \underline{69.4\%} & 25.0\% & \underline{61.5\%} & \underline{77.3\%} \\
\bottomrule
\end{tabular}%
}
\end{table*}

\subsection{Verifiable Problems}

% Verifiable tasks provide the cleanest test of whether the post-training recipe
%   improves problem solving under trusted evaluation. 
  As shown in
  \cref{tab:verifiable-single-pass}, SU-01 reaches a 77.3\% average score across
  AnswerBench, AMO-Bench, AIME 2025, AIME 2026, and FrontierScience-Olympiad,
  nearly matching the strongest similar-size baseline, Qwen3.6-35B-A3B (77.4\%).
  Importantly, SU-01 achieves this level of performance with a simpler unified
  post-training recipe and substantially lower training cost, highlighting the
  efficiency of our approach.
The mathematical benchmarks show where this improvement is most pronounced.
SU-01 achieves the best similar-size results on AMO-Bench (59.8\%) and AIME
2025/2026 (94.6\%/93.3\%), which are closer to competition-style problem solving
than routine answer extraction. On AnswerBench, SU-01 remains competitive at
77.5\%, close to Qwen3.6-35B-A3B (78.0\%) and behind only Nemotron-Cascade-2
(80.5\%). 
% This pattern suggests that the recipe is especially effective at
% strengthening hard mathematical reasoning while preserving broad answer-verifiable
% competence.
  FrontierScience-Olympiad tests whether this behavior transfers beyond pure
  mathematics. Although the RL stages use only math and physics signals, SU-01
  reaches 61.5\% overall and shows strong transfer to untrained STEM domains,
  including 69.4\% on Chemistry and 25.0\% on Biology. This cross-domain transfer
  supports the specializable-generalist framing: the model is specialized through
  math and physics reasoning signals, yet the resulting capability does not
  collapse into a narrow contest solver.
%  rather than simply
% memorizing the reward domains used during RL.



\subsection{Non-verifiable Problems}

Non-verifiable benchmarks test whether the training recipe improves the quality
of full reasoning traces, rather than only optimizing final-answer rewards. On
IMO-ProofBench, \cref{tab:nonverifiable-benchmarks} shows that SU-01 reaches
57.6\% overall in direct generation, already the strongest result among
similar-size models. Test-time scaling further raises the score to 70.2\%,
including 91.0\% on the basic split and 49.5\% on the advanced split, bringing a
30B-A3B model close to much larger frontier systems such as Gemini 3.1 Pro. This improvement indicates
that self-verification and refinement are especially useful when correctness
depends on the complete proof, not merely on producing the right final answer.

\begin{wraptable}{r}{0.54\textwidth}
  % \vspace{-1.1em}
  \centering
  \captionsetup{type=table,hypcap=false}
  \caption{Performance on physics olympiad problems, averaged over 8 runs. Gold lines for IPhO 2024/2025 are
  20.8/19.7 points. For SU-01, $x/y$ reports scores without and with test-time scaling. 
  % Within each block,
  % \textbf{bold} marks the best score and underline marks the second best.
  }
  \vspace{-0.75em}
  \label{tab:ipho-problems}
  \setlength{\tabcolsep}{3.0pt}
  \renewcommand{\arraystretch}{1.}
  \small
  \resizebox{0.85\linewidth}{!}{%
  \begin{tabular}{@{}lc@{\hspace{1.0em}}c@{}}
  \toprule
  \textbf{Model} & \textbf{IPhO 2024} & \textbf{IPhO 2025} \\
  \midrule
  \multicolumn{3}{l}{\textit{Larger models}} \\
  \midrule
  Gemini 3.1 Pro Thinking & \textbf{25.9} & \textbf{25.1} \\
  GPT-5.5-High & \underline{25.8} & \underline{23.2} \\
  DeepSeek-V3.2-Speciale & 25.1 & 21.9 \\
  \midrule
  \multicolumn{3}{l}{\textit{Similar-size models}} \\
  \midrule
  P1-30B-A3B & 23.1 & 17.7 \\
  GLM-4.7-Flash & 22.2 & 19.5 \\
  Nemotron-Cascade-2 & 21.2 & 16.7 \\
  Qwen3.6-35B-A3B & 24.3 & 19.9 \\
  Gemma-4-31B & \underline{24.4} & \underline{20.3} \\
  \rowcolor{ScienceReasoningBg!35}
  \textbf{\textcolor{iclrdeepblue}{SU-01}} & 23.5\,/\,\textbf{25.3} & \underline{20.3}\,/\,\textbf{21.7} \\
  \bottomrule
  \end{tabular}
  }
  \vspace{-1.0em}
  \end{wraptable}
  
FrontierScience-Research is a substantially harder research-oriented subset of
FrontierScience, covering physics, chemistry, and biology problems that require
scientific modeling and multi-step reasoning beyond standard contest formats.
Absolute scores remain low even for frontier systems, but SU-01 obtains the
best similar-size overall score at 11.7\%. It also leads the
similar-size block on Physics, ties for the best Chemistry score, and ranks
second on Biology, despite our RL stages using only mathematics and physics
signals. This cross-domain pattern suggests that the recipe learns a more
general scientific reasoning behavior rather than only specializing to the
training domains, providing early evidence of transferable research-level
reasoning in a compact model.

\begin{table*}[t]
\centering
\caption{Performance on non-verifiable benchmarks. FrontierScience-Research refers to the research subset of FrontierScience. For SU-01, $x/y$ reports scores without and with TTS on ProofBench.}
\label{tab:nonverifiable-benchmarks}
% \begin{threeparttable}
\vspace{-5 pt}
\setlength{\tabcolsep}{2.4pt}
\renewcommand{\arraystretch}{1.12}
\small
\resizebox{0.97\textwidth}{!}{%
\begin{tabular}{lccccccc}
\toprule
\multirow{2}{*}{\textbf{Model}} & \multicolumn{3}{c}{\textbf{IMO-ProofBench}} & \multicolumn{4}{c}{\textbf{FrontierScience-Research}} \\
\cmidrule(lr){2-4}\cmidrule(lr){5-8}
 & \textbf{Basic} & \textbf{Advanced} & \textbf{Overall} & \textbf{Physics} & \textbf{Chemistry} & \textbf{Biology} & \textbf{Overall} \\
\midrule
\multicolumn{8}{l}{\textit{Larger models}} \\
\midrule
Gemini 3.1 Pro Thinking & \underline{95.2\%} & \underline{50.0\%} & \underline{72.6\%} & 0.0\% & \underline{30.0\%} & 10.0\% & 13.3\% \\
GPT-5.5-High & \textbf{96.7\%} & \textbf{64.8\%} & \textbf{80.7\%} & \textbf{25.0\%} & \textbf{40.0\%} & \textbf{45.0\%} & \textbf{36.7\%} \\
DeepSeek-V3.2-Speciale & 77.6\% & 34.3\% & 56.0\% & \underline{10.0\%} & 20.0\% & \underline{15.0\%} & \underline{15.0\%} \\
\midrule
\multicolumn{8}{l}{\textit{Similar-size models}} \\
\midrule
P1-30B-A3B & 33.8\% & 6.2\% & 20.0\% & 0.0\% & \textbf{10.0\%} & 0.0\% & 3.3\% \\
GLM-4.7-Flash & 51.0\% & 16.7\% & 33.8\% & 0.0\% & 0.0\% & 0.0\% & 0.0\% \\
Nemotron-Cascade-2 & \underline{77.1\%} & 28.6\% & 52.9\% & \underline{5.0\%} & 5.0\% & \textbf{20.0\%} & \underline{10.0\%} \\
Qwen3.6-35B-A3B & 39.1\% & 7.1\% & 23.1\% & 0.0\% & 5.0\% & 10.0\% & 5.0\% \\
Gemma-4-31B & 46.7\% & 16.2\% & 31.4\% & 0.0\% & \textbf{10.0\%} & 5.0\% & 5.0\% \\
\rowcolor{ScienceReasoningBg!35}
\textbf{\textcolor{iclrdeepblue}{SU-01}} & \underline{77.1\%}\,/\,\textbf{91.0\%} & \underline{38.1\%}\,/\,\textbf{49.5\%} & \underline{57.6\%}\,/\,\textbf{70.2\%} & \textbf{10.0\%} & \textbf{10.0\%} & \underline{15.0\%} & \textbf{11.7\%} \\
\bottomrule
\end{tabular}%
}
% \begin{tablenotes}
% \footnotesize
% \item[*]Please refer to Appendix \ref{app:inference-serving-details} for the evaluation details of DeepSeek-V3.2-Speciale.
% \end{tablenotes}
% \end{threeparttable}
\end{table*}

% \clearpage
\subsection{Olympiad Competition Problems}

% Olympiad problems provide a direct view of whether the model can translate its
% reasoning improvements into competition-level performance under strict scoring
% criteria. We separate these results from the general verifiable benchmarks
% because they are closer to the final target setting of the report.


\begin{wraptable}{r}{0.58\textwidth}
  \vspace{-1.3em}
  \centering
  \captionsetup{type=table,hypcap=false}
  \caption{Performance on mathematical olympiad competition problems. Medal lines for
  IMO 2025 are 35/28/19 points, and medal lines for USAMO 2026 are 25/18/11 points.
  % $^{\star}$ marks human-expert evaluation for TTS scores; direct rows use
  % automatic evaluation under the same protocol (\cref{app:evaluation-details}).
  $^{\star}$ indicates that TTS results are evaluated by human experts, while direct generation results are evaluated automatically (\cref{app:evaluation-details}).
  }
  \vspace{-0.75em}
  \label{tab:olympiad-competition-problems}
  \setlength{\tabcolsep}{5.8pt}
  \renewcommand{\arraystretch}{1.0}
  \small
  \begin{tabular}{@{}lccccccc@{}}
  \toprule
  \multicolumn{8}{c}{\textbf{IMO 2025}} \\
  \midrule
  \textbf{Model} & \textbf{P1} & \textbf{P2} & \textbf{P3} & \textbf{P4} & \textbf{P5} & \textbf{P6} & \textbf{Total} \\
  \midrule
  % \rowcolor{MathReasoningBg!16}
  \textbf{\textcolor{iclrdeepblue}{SU-01}} & 1 & 7 & 1 & 6 & 6 & 0 & 21 \\
  % \rowcolor{MathReasoningBg!16}
  \textbf{\textcolor{iclrdeepblue}{SU-01} w/ TTS} & $7^{\star}$ & $7^{\star}$ & $7^{\star}$ & $7^{\star}$ & $7^{\star}$ & $0^{\star}$ & $35^{\star}$\,\goldmedalicon \\
  \midrule
  \multicolumn{8}{c}{\textbf{USAMO 2026}} \\
  \midrule
  \textbf{Model} & \textbf{P1} & \textbf{P2} & \textbf{P3} & \textbf{P4} & \textbf{P5} & \textbf{P6} & \textbf{Total} \\
  \midrule
  % \rowcolor{MathReasoningBg!16}
  \textbf{\textcolor{iclrdeepblue}{SU-01}} & 7 & 0 & 0 & 7 & 0 & 1 & 15 \\
  % \rowcolor{MathReasoningBg!16}
  \textbf{\textcolor{iclrdeepblue}{SU-01} w/ TTS} & $7^{\star}$ & $0^{\star}$ & $7^{\star}$ & $7^{\star}$ & $7^{\star}$ & $7^{\star}$ & $35^{\star}$\,\goldmedalicon \\
  \bottomrule
  \end{tabular}
  % \vspace{-2.0em}
  \end{wraptable}



% \cref{tab:ipho-problems} shows strong transfer to physics olympiad settings.
Even without TTS, SU-01 averages 23.5 points on IPhO 2024 and 20.3 points on
IPhO 2025, exceeding the corresponding gold lines of 20.8 and 19.7 points, as shown in \cref{tab:ipho-problems}. TTS
further raises the scores to 25.3 and 21.7 points, making SU-01 the strongest
similar-size model in both years among models with available scores. 
% This is a
% notable cross-domain result: although the training recipe is centered on
% rigorous reasoning and math/physics feedback, the model must still combine
% physical modeling, derivation, numerical reasoning, and competition-style
% justification to score well on IPhO problems.

\cref{tab:olympiad-competition-problems} reports the final competition-style
mathematics results. In direct generation, SU-01 reaches 21 points on IMO 2025
and 15 points on USAMO 2026, already clearing the bronze-medal lines for both
competitions. 
% The problem-level pattern is uneven but informative: 
The direct
model obtains full credit on IMO 2025 P2 and USAMO 2026 P1/P4, near-complete
solutions on IMO 2025 P4/P5, and still fails several harder problems in a
single pass. 
This indicates that the base model has acquired substantial
olympiad reasoning ability, but still benefits from additional search and
self-correction on the most brittle proof attempts.

With test-time scaling, SU-01 reaches \textbf{35 points} on both IMO 2025 and USAMO 2026, meeting the \textbf{IMO gold} line exactly and exceeding the \textbf{USAMO gold} line by 10
points. TTS upgrades five IMO 2025 problems to full credit while P6 remains unsolved, and recovers full-credit solutions on five of six USAMO 2026 problems while
P2 remains unresolved. The USAMO 2026 score summary reports 340 competitors, a median score of 6, a top-12 cutoff of 26, and a maximum score of 35\footnote{https://web.evanchen.cc/exams/posted-usamo-statistics.pdf}. 
SU-01
therefore matches the \textbf{highest} reported human total on this contest, while still exposing a concrete failure mode on P2. 
This result suggests that our overall
recipe can elicit top-level human-like olympiad reasoning from a compact 30B-A3B model.


\paragraph{Case study.} 
We include the corresponding
model-generated solutions and expert verdicts in
\cref{app:olympiad-solutions}. 
Across the twelve IMO 2025 and USAMO 2026 problems, the model gives full-credit solutions to ten problems, with failures on IMO 2025 P6 and USAMO 2026 P2. Its main strength is translating olympiad problems into formal frameworks: coordinates or complex numbers for geometry, modular classifications for number theory, recurrences for functional equations, and automata-based dynamic programming for digit problems. A particularly striking example is USAMO 2026 P3: rather than following the standard synthetic geometry route, the model elegantly uses complex numbers to unify the unit circle, equilateral-triangle rotations, chord relations, and tangent conditions within a single algebraic framework. This yields an ingenious analytic reformulation of a configuration that olympiad solvers would typically approach through angle chasing and carefully chosen auxiliary constructions. IMO 2025 P2 shows a complementary strength: the model reduces a configuration involving two intersecting circles, an orthocenter, and a tangency claim to coordinate and distance computations. Other strong examples include the carry-state dynamic programming approach for USAMO P4 and the number-theoretic proof using totients, congruences, Vieta jumping, and Fibonacci structure in USAMO P6. However, the failures show a clear limitation: the model can miss subtle structural constraints, as in the invalid column-permutation reduction in IMO P6, or leave gaps in delicate global strategy arguments, as in USAMO P2. Overall, the model performs well when a problem admits a rigid formal representation, but is less reliable when the core challenge is preserving combinatorial structure or proving a finely tuned process invariant.